\documentclass[conference]{IEEEtran}
\usepackage{amsmath,amssymb,booktabs,graphicx,xcolor}
\usepackage[hidelinks]{hyperref}
\newcommand{\E}{\mathbb E}
\newcommand{\pending}{\textemdash}
\newcommand{\draftnote}[1]{\textcolor{blue}{[Draft: #1]}}
% Six-page target, including references; final pagination requires compilation.
% Approximate page allocation: introduction 0.75; formulation 1.25;
% methods 1.25; setup 0.75; results/discussion 1.5; conclusion/references 0.5.
% Populate pending results only from completed, verified conference campaigns.
% Full project reference: objective_driven_boed_manuscript_framework.tex.
\title{Sequential Bayesian Experimental Design for Non-Reset
Power-Grid Parameter Learning}
\author{\IEEEauthorblockN{Authors to be supplied}
\IEEEauthorblockA{Affiliations to be supplied}}
\begin{document}
\maketitle
\begin{abstract}
Controlled probing can reveal uncertain power-system parameters, but each
probe also changes the state entering the next experiment. We study sequential
Bayesian experimental design for this non-reset setting using reduced IEEE9
and IEEE14 models with uncertain inertia and effective frequency response.
The design variable is a continuous probe duration; observations are noisy
signed endpoint rate-of-change-of-frequency summaries. We compare DAD,
RL-sBOED, Step-DAD, Myopic, Fixed and Random under common physical constraints
and paired evaluation systems. The study evaluates a sequential
prior-contrastive lower bound on expected information gain and separates
offline training from online design cost. Step-DAD uses four conditional
refinement updates after each observation. The protocol covers three horizons
and three independent training seeds. Numerical assumptions and the
finite-contrast ceiling are stated explicitly.
\draftnote{Add the principal quantitative finding after the campaign is complete.}
\end{abstract}
\begin{IEEEkeywords}
Bayesian experimental design, power-system identification, adaptive probing,
expected information gain, rate of change of frequency.
\end{IEEEkeywords}

\section{Introduction}
Inertia and effective frequency response govern how a power system reacts
to disturbances. Designed excitation can provide information about both,
motivating probing-based parameter estimation~\cite{peng2024probing}.
Here the question is how to choose successive probe durations when previous
observations inform the next choice and the physical system does not reset.
The reduced dynamics and scalar sensor used below are study assumptions;
they are not a reproduction of a hardware-validated identification system.

A sequential design must account for both information and physical state.
An early action determines its observation and changes the state from which
later actions are applied. Replacing those later states with equilibrium
would change the experiment. Likewise, histories must preserve observation
order. These requirements motivate comparing full-horizon policies with
one-step optimization and observation-independent baselines.

DAD amortizes design optimization into an offline-trained policy and supplies
sequential contrastive information bounds~\cite{foster2021dad}.
RL-based design learns long-horizon behavior through simulated rewards
without simulator derivatives~\cite{blau2022rlsboed}. Step-DAD refines an
offline policy after observations arrive~\cite{hedman2025stepdad}.
We adapt these families to ordered, state-carrying power-grid experiments.
Our contribution is the common application formulation and comparison,
including the information--online-computation trade-off, rather than a claim
that these policy families are new or that adaptivity guarantees improvement.

The paper addresses three questions: whether observation-dependent policies
improve information over Fixed and Random; whether full-horizon policies
improve on Myopic; and whether online refinement provides sufficient benefit
to justify its computational cost. The scope is EIG on IEEE9/IEEE14.
MSC/MOCU control objectives, SIR, and optional MoE-sBOED remain in the full
project document and are not additional experiments in this comparison.

\section{Non-Reset Experiment and Information Objective}
\subsection{Reduced Dynamics and Parameter Uncertainty}
IEEE9 retains three dynamic machines and IEEE14 retains five. Other buses
are algebraic network nodes, not independent inertial machines. The network
reduction and injection map use the declared test-case data
\cite{zimmerman2011matpower}. For retained machine $i$,
\begin{align}
\dot\delta_i &= \omega_i,\\
M_i\dot\omega_i &= -\sum_j B_{ij}
 [\sin(\delta_i-\delta_j)-\sin(\delta_i^0-\delta_j^0)]
 -(D_i+K_i/(2\pi))\omega_i+v_i .
\label{eq:swing}
\end{align}
Here $\omega_i$ is angular-speed deviation, $\Delta f_i=\omega_i/(2\pi)$,
$D_i$ is known damping, and $K_i$ is effective response in per-unit power
per Hz. The uncertain vector is
$\theta=(M_1,\ldots,M_m,K_1,\ldots,K_m)$, with independent bounded uniform
priors fixed before training. Each system retains its machine-specific
prior bounds, damping and coupling; parameter values are not transferred
between IEEE9 and IEEE14. The reduced sinusoidal model with a linear Kron
injection map is an approximation, not a full AC transient-stability model.
\draftnote{Provide the final configuration archive/reference containing exact prior bounds and network parameters.}

\subsection{Probe Sequence and Observation}
A probe at fixed physical bus 1 has fixed amplitude $A=0.05$ pu and duration
$\tau_t$. For local time $\ell$ within a window of length $W$,
\begin{equation}
\mathbf v_t(\ell)=\mathbf a(1)\frac{A}{2}
 [1-\cos(2\pi\ell/\tau_t)]\mathbf1_{[0,\tau_t]}(\ell).
\end{equation}
Writing $\mathbf x=(\boldsymbol\delta,\boldsymbol\omega)$, the transition is
$\mathbf x_t=\Phi_\theta(\mathbf x_{t-1},\tau_t;W)$.
Only episode initialization uses equilibrium; all subsequent machine states
carry forward. The same latent parameters apply throughout the episode.

Durations lie in $[0.2,3]$ s and satisfy increasing order with separation
$\varepsilon=0.01$ s. Before stage $t$, define
\begin{equation}
a_t=\begin{cases}0.2,&t=1,\\\tau_{t-1}+\varepsilon,&t>1,\end{cases}
\quad b_t=3-(T-t)\varepsilon.
\end{equation}
Every method maps a unit action $v_t\in[0,1]$ to
$\tau_t=a_t+(b_t-a_t)v_t$ before simulation. This reserves space for later
probes; no completed history is sorted or rounded to a duration grid.

The observation at bus 1 is signed endpoint RoCoF:
\begin{equation}
y_t=\frac{\Delta f_1(s_t+W)-\Delta f_1(s_t+W-\Delta r)}{\Delta r}
 +\epsilon_t,
\label{eq:sensor}
\end{equation}
where $s_t=(t-1)W$, $W=3.5$ s, $\Delta r=0.025$ s, and
$\epsilon_t\sim\mathcal N(0,0.005^2)$ independently across windows, in
Hz/s. Noise is added to the clean summary. This is not maximum absolute
RoCoF. Observations arrive at window end. Computation latency is zero in
simulated time; measured decision runtime is reported separately.

\subsection{Sequential Prior-Contrastive Information}
For ordered history $h_t=((\tau_k,y_k))_{k=1}^t$, a policy chooses
$\tau_t=\pi(h_{t-1})$. The information objective is
\begin{equation}
J(\pi)=\E_{\theta,H_T\mid\pi}
\left[\log\frac{p(\theta\mid H_T)}{p_0(\theta)}\right].
\end{equation}
We optimize and report the sequential prior-contrastive (sPCE) estimator
\cite{foster2021dad}. One model $\theta^{(0)}$ generates the observations;
$L$ independent models are drawn from the prior. Define
\begin{align}
\ell_t^{(n)}&=\sum_{k=1}^t\log p(y_k\mid\theta^{(n)},h_{k-1},\tau_k),\\
\widehat I_t&=\ell_t^{(0)}-
\log\left[\frac{1}{L+1}\sum_{n=0}^L\exp\ell_t^{(n)}\right].
\label{eq:spce}
\end{align}
Each model follows the actual selected durations while propagating its own
state. The expected estimator is a lower bound on EIG; an individual score
can be negative. Its upper ceiling is $\log(L+1)$, equal to 4.8598 nats for
the conference setting $L=128$. Near-ceiling scores cannot establish that
true EIG has saturated or resolve differences hidden by the estimator.
Evaluation truth appears in this scoring denominator, never as a privileged
input to an online planner.

\section{Compared Design Methods}
\subsection{Common History and Offline Policies}
Policy inputs contain the stage index and ordered slots for past durations,
scaled observations and availability masks. Future slots are masked.
True parameters and simulator states are not policy inputs.

\textbf{DAD} uses a deterministic two-layer, 64-unit tanh network with a
trainable stage bias and sigmoid output. Full-horizon sPCE training propagates
gradients through actions, observations and carried states. Joint Fixed/DAD
training selects between random and interior Fixed-based initialization
using validation only; the Fixed initialization cost must be disclosed.
The selected DAD policy requires no online weight updates.

\textbf{RL-sBOED} uses a REDQ-style continuous actor--critic adaptation with
ten critics and five replay updates per trajectory batch. Prefix rewards
$r_t=\widehat I_t-\widehat I_{t-1}$ telescope to terminal sPCE with
$\gamma=1$. Thus stepwise rewards do not make it myopic. The actor has two
128-unit ReLU layers. Entropy uses the normalized tanh action with target
entropy $-1$; evaluation scores exclude entropy bonuses. This adaptation
does not require derivatives through the simulator.

\subsection{Step-DAD Online Refinement}
Step-DAD starts each episode from the matching DAD checkpoint. The first
action uses that policy unchanged. Following each observation, 128 planning
particles update likelihood weights and carry their own states under the
observed action sequence. Fantasy truths and contrasts are sampled from
this posterior to optimize the remaining-horizon conditional objective.

Each refinement performs \textbf{four gradient updates}, using batch size
32 and learning rate $10^{-3}$. A separate fixed conditional validation
batch compares the initial and updated policies; the best candidate is
retained. Adapted weights carry to the next stage, while the optimizer is
reinitialized at each refinement. Episodes do not share adapted weights.
There are $4(T-1)$ updates per episode: 8, 12 and 16 at the three horizons.
The earlier 16-updates-per-stage results are superseded and excluded.

\subsection{Baselines}
\textbf{Fixed} optimizes one ordered, observation-independent sequence
offline, comparing interior short, central, long and spread initializations.
\textbf{Myopic} maximizes expected one-step information through continuous
search with 16 proposals, three rounds and 32 fantasies per proposal.
\textbf{Random} samples uniformly from each remaining feasible interval;
this is not uniform sampling over complete ordered sequences.

\subsection{Numerical Verification}
RK4 integration uses step $1/640$ s. Numerical simulator Jacobians use
state and duration perturbations of order $10^{-5}$, including the
carried-state chain rule. Full-policy and conditional-prefix gradients
must agree with independent finite-difference checks. The endpoint sensor
uses fixed sample indices and does not require peak-switch differentiation.
Finite-particle inference, ODE discretization and policy optimization are
distinct error sources. Methods are documented adaptations, not claims of
exact reproduction of their original benchmarks.

\section{Experimental Protocol}
\begin{table}[t]
\caption{Matched conference settings for IEEE9 and IEEE14.}
\centering\small
\begin{tabular}{ll}
\toprule Setting & Value\\\midrule
Horizons & $T=3,4,5$\\
Training seeds & 101, 202, 303\\
Offline updates / batch size & 2,000 / 32\\
DAD/Fixed learning rate & $10^{-3}$\\
Validation & 128 systems, every 100 updates\\
sPCE contrasts / planner particles & 128 / 128\\
Evaluation seeds & 1001, 1002, 1003\\
Systems per evaluation seed & 512\\
Step-DAD updates after observation & 4\\
\bottomrule
\end{tabular}
\label{tab:settings}
\end{table}
IEEE9 and IEEE14 use the same sensor, feasible design rule and numerical
budgets, but separate physical configurations and learned policies.
Validation selects checkpoints without gradient updates on validation cases.
Each trained policy is evaluated on 1,536 episodes. Parameter and
measurement-noise draws are paired across methods; the online planner uses
a separate random stream. Random/Myopic have no training-seed dependence
under these streams, so identical evaluations can be shared across training
seeds. Fixed remains independently optimized for each training seed.

For each trained method, average the three evaluation-seed means within a
training seed, then report mean $\pm$ sample standard deviation across
the three training seeds. Nine training/evaluation combinations do not
constitute nine independent trainings. System dispersion and paired
within-policy uncertainty are separate from retraining variation.
Seed 1001 has informed development; it is not untouched confirmatory data.

Record offline optimization time and online design time separately.
Online episode time is the sum of decision times over all stages and
includes posterior updates/refinement where used. Report CPU/GPU hardware
alongside timings. IEEE9 uses Ubuntu RTX 5080 and IEEE14 uses Grace A100;
cross-system timing differences therefore cannot be attributed solely to
network size. EIG scoring time is not interchangeable with decision latency.

\section{Results and Discussion}
\draftnote{This section awaits final verified aggregation. Dashes below mean
pending, not zero. Do not insert pilot scores or superseded Step-DAD results.}
\subsection{Information Across Horizons}
\begin{table}[t]
\caption{Terminal sPCE (nats), IEEE9. Trained methods: mean $\pm$ sample
SD across three training seeds. $L=128$; ceiling 4.8598. Results pending.}
\centering
\begin{tabular}{lccc}
\toprule Method & $T=3$ & $T=4$ & $T=5$\\\midrule
Random & \pending & \pending & \pending\\
Fixed & \pending & \pending & \pending\\
Myopic & \pending & \pending & \pending\\
DAD & \pending & \pending & \pending\\
RL-sBOED & \pending & \pending & \pending\\
Step-DAD & \pending & \pending & \pending\\\bottomrule
\end{tabular}
\label{tab:ieee9}
\end{table}
\begin{table}[t]
\caption{Terminal sPCE (nats), IEEE14, using the same aggregation and
contrast count as Table~\ref{tab:ieee9}. Results pending.}
\centering
\begin{tabular}{lccc}
\toprule Method & $T=3$ & $T=4$ & $T=5$\\\midrule
Random & \pending & \pending & \pending\\
Fixed & \pending & \pending & \pending\\
Myopic & \pending & \pending & \pending\\
DAD & \pending & \pending & \pending\\
RL-sBOED & \pending & \pending & \pending\\
Step-DAD & \pending & \pending & \pending\\\bottomrule
\end{tabular}
\label{tab:ieee14}
\end{table}
The primary comparisons are learned policies versus Fixed/Random,
full-horizon policies versus Myopic, and Step-DAD versus its matching DAD.
Report paired differences rather than interpreting small mean differences
alone as evidence of a ranking. Random/Myopic do not have a retraining SD;
their uncertainty must be labeled separately. Near-ceiling IEEE14 scores
require particular care when interpreting apparent ties.

\subsection{Information Versus Online Computation}
\begin{figure}[t]
\centering
\fbox{\parbox[c][1.1in][c]{0.91\columnwidth}{\centering
Pending verified figure:\\sPCE versus online seconds per episode\\
logarithmic time axis; separate IEEE9/IEEE14 panels}}
\caption{Planned information--computation comparison. Use measured
online time on the stated hardware and label horizons. No data are drawn yet.}
\label{fig:tradeoff}
\end{figure}
Compare the information benefit of refinement with its additional online
cost, distinguishing offline-trained feed-forward policies from online
simulation-based planners. A fast simulated decision does not certify
feasibility for a physical deadline, particularly because the experiment
model assumes zero decision delay.

\subsection{Convergence and Interpretation Limits}
Include a compact validation-history panel for representative completed
training runs, showing initialization and selected checkpoints. A fixed
2,000-update budget is not evidence of convergence. Contrast sensitivity,
if available, should evaluate frozen policies on paired cases; a score below
the ceiling alone does not establish estimator accuracy. Particle sensitivity
addresses online inference and is a different issue from contrast count.
Only completed diagnostics may be reported as evidence.

The reduced dynamics, selected priors and scalar sensor limit generalization.
The paper reports information, not terminal-control safety or field validation.
Optional MoE-sBOED and control-oriented objectives are reserved for the
broader project; no MoE superiority claim is part of this six-method study.

\section{Conclusion}
This study specifies a common non-reset power-grid experiment for comparing
amortized, semi-amortized and online Bayesian design. The evaluation separates
finite-contrast information from online computational cost under matched
physical assumptions. \draftnote{Replace this sentence with conclusions
supported by the completed IEEE9/IEEE14 tables and timing figure.}
\bibliographystyle{IEEEtran}
\bibliography{objective_driven_boed_refs}
\end{document}
